\documentclass[12pt]{article}
\usepackage[utf8]{inputenc}
\usepackage[T1]{fontenc}
\usepackage{hyperref}
\title{Taman Language Resources for K-12 Curriculum Development}
\author{Compiled by Claude Sonnet 5 (claude-sonnet-5) for Andr\'as Kornai}
\date{}
\begin{document}
\maketitle

\section{Introduction}
Taman (ISO 639-3: tmn) is an Austronesian language of Borneo, spoken mainly in Kapuas Hulu Regency, West Kalimantan, Indonesia. With Embaloh it forms the Tamanic branch, which is related to the languages of South Sulawesi rather than to other Bornean languages. One source gives about 30,000 speakers (2007). Do not confuse it with Embaloh, a separate Tamanic language. Search results showed little that is specific to Taman.

\section{Primary Educational Resources}
No primer, textbook or school programme for Taman turned up in search results.

\section{Linguistic and Reference Materials}
Wikipedia has an article on the language: \url{https://en.wikipedia.org/wiki/Taman_language_%28Indonesia%29}, and one on the Tamanic branch: \url{https://en.wikipedia.org/wiki/Tamanic_languages}. The Embaloh article is related but covers a different language: \url{https://en.wikipedia.org/wiki/Embaloh_language}. A dialectology article on Tamanic appeared in the results, but I did not open it, so its content and relevance are unconfirmed: \url{https://raco.cat/index.php/Dialectologia/article/download/427266/521774/628251}. No Taman dictionary, grammar or orthography was found.

\section{Religious and Community Resources}
None found for Taman. The Joshua Project and Sabda results were for Embaloh, a different language, so they are not listed here.

\section{Digital and Audio Resources}
None found.

\section{Curriculum Development Context}
The language is not known to be taught in schools, and the number of young speakers is reportedly decreasing. A curriculum would have to start from documentation work with the community.
\end{document}